\begin{lemma}[Mollifications of an $L^2_{\mathrm{loc}}$ function converge in $L^2$ on compact sets]
Write $\lambda$ for Lebesgue measure on $\mathbb{R}^n$ and let
$f : \mathbb{R}^n \to \mathbb{R}$ be such that $f \in L^2(K,\lambda)$ for every
compact $K \subseteq \mathbb{R}^n$; that is, $f \in L^2_{\mathrm{loc}}(\mathbb{R}^n)$.
Let $(\rho_\varepsilon)_{\varepsilon>0}$ be a mollifier family: for each
$\varepsilon > 0$ the function $\rho_\varepsilon : \mathbb{R}^n \to \mathbb{R}$ is
smooth, satisfies $\rho_\varepsilon \geq 0$, is supported in the closed ball
$\{|t| \leq \varepsilon\}$, and has $\int_{\mathbb{R}^n}\rho_\varepsilon\,d\lambda = 1$.
Write
\[
  (\rho_\varepsilon * f)(x) = \int_{\mathbb{R}^n} \rho_\varepsilon(t)\, f(x-t)\, d\lambda(t) .
\]
Then for every compact $K \subseteq \mathbb{R}^n$,
\[
  \bigl\|\rho_\varepsilon * f - f\bigr\|_{L^2(K,\lambda)} \longrightarrow 0
  \qquad\text{as } \varepsilon \downarrow 0
\]
through positive values of $\varepsilon$, the norms being read in $[0,+\infty]$.
No continuity, boundedness or global integrability of $f$ is assumed.
\end{lemma}
\begin{proof}
Fix the compact set $K$.  For $r > 0$ put
$K_r = \{x \in \mathbb{R}^n : \operatorname{dist}(x,K) \leq r\}$; each $K_r$ is
closed and bounded, hence compact, and $\lambda(K) \leq \lambda(K_r) < \infty$.
Throughout, $\varepsilon$ is restricted to $0 < \varepsilon \leq 1$.  If
$K = \emptyset$ every norm below is $0$ and there is nothing to prove, so assume
$K \neq \emptyset$.

\emph{Step 0: the convolution is defined on $K$.}
Let $0 < \varepsilon \leq 1$ and $x \in K$.  Since
$\operatorname{supp}\rho_\varepsilon \subseteq \{|t|\leq\varepsilon\}$, the
integrand $t \mapsto \rho_\varepsilon(t) f(x-t)$ vanishes unless
$|t| \leq \varepsilon \leq 1$, and for such $t$ we have $x - t \in K_1$.  Now
$f \in L^2(K_1,\lambda)$ by hypothesis and $\lambda(K_1) < \infty$, so
Cauchy--Schwarz (shared context) on the measure space $(K_1,\lambda)$, applied to
$|f|$ and the constant $1$, gives
$\int_{K_1}|f|\,d\lambda \leq \|f\|_{L^2(K_1,\lambda)}\lambda(K_1)^{1/2} < \infty$.
As $\rho_\varepsilon$ is continuous with compact support it is bounded, say by
$M_\varepsilon$, so
$\int |\rho_\varepsilon(t) f(x-t)|\,d\lambda(t) \leq M_\varepsilon \int_{K_1}|f|\,d\lambda < \infty$
after the change of variable $y = x-t$, which preserves $\lambda$ by translation
invariance and reflection invariance of Lebesgue measure (shared context).  Hence
$(\rho_\varepsilon * f)(x)$ is a well-defined real number for every $x \in K$ and
every $0 < \varepsilon \leq 1$.

\emph{Step 1: a uniform local $L^2$ bound for the mollification.}
\emph{Claim.}  For every measurable $h : \mathbb{R}^n \to \mathbb{R}$ with
$h \in L^2(K_1,\lambda)$ and every $0 < \varepsilon \leq 1$,
\begin{equation*}
  \|\rho_\varepsilon * h\|_{L^2(K,\lambda)} \;\leq\; \|h\|_{L^2(K_1,\lambda)} .
  \tag{1}
\end{equation*}
\emph{Proof of the claim.}  Fix $0 < \varepsilon \leq 1$.  Because
$\rho_\varepsilon \geq 0$ and $\int \rho_\varepsilon \, d\lambda = 1$, the measure
$\mu_\varepsilon = \rho_\varepsilon\,d\lambda$ is a probability measure on
$\mathbb{R}^n$.  Fix $x \in K$.  The function
$t \mapsto h(x-t)$ is $\mu_\varepsilon$-integrable by the computation of Step 0
applied to $h$, and $s \mapsto s^2$ is convex on $\mathbb{R}$, so Jensen's
inequality for a probability measure (shared context) gives
\[
  \bigl((\rho_\varepsilon * h)(x)\bigr)^2
    = \Bigl(\int h(x-t)\, d\mu_\varepsilon(t)\Bigr)^{2}
    \;\leq\; \int h(x-t)^2 \, d\mu_\varepsilon(t)
    \;=\; \int \rho_\varepsilon(t)\, h(x-t)^2\, d\lambda(t) .
\]
The function $(x,t) \mapsto \rho_\varepsilon(t) h(x-t)^2$ is nonnegative and
jointly measurable, being the product of a measurable function of $t$ with the
composition of the measurable $h^2$ with the continuous map $(x,t)\mapsto x-t$.
Integrating the displayed inequality over $x \in K$ and applying Tonelli's
theorem (shared context) to interchange the order of integration,
\[
  \int_K \bigl((\rho_\varepsilon * h)(x)\bigr)^2 d\lambda(x)
    \;\leq\; \int \rho_\varepsilon(t) \Bigl(\int_K h(x-t)^2\, d\lambda(x)\Bigr) d\lambda(t) .
\]
For $|t| \leq \varepsilon \leq 1$, translation invariance of Lebesgue measure
(shared context) gives
$\int_K h(x-t)^2 d\lambda(x) = \int_{K-t} h(y)^2 d\lambda(y)$, and
$K - t \subseteq K_1$ because every $y \in K-t$ satisfies
$\operatorname{dist}(y,K) \leq |t| \leq 1$; since $h^2 \geq 0$ this yields
$\int_K h(x-t)^2 d\lambda(x) \leq \|h\|_{L^2(K_1,\lambda)}^2$.  For $|t| > \varepsilon$
the factor $\rho_\varepsilon(t)$ vanishes.  Hence the right-hand side is at most
$\|h\|_{L^2(K_1,\lambda)}^2 \int \rho_\varepsilon\, d\lambda = \|h\|_{L^2(K_1,\lambda)}^2$.
Taking square roots gives $(1)$ and proves the claim.

\emph{Step 2: the case of a continuous, compactly supported function.}
Let $g : \mathbb{R}^n \to \mathbb{R}$ be continuous with compact support.  By the
density item of shared context, such a $g$ is bounded and uniformly continuous;
let $\omega(\varepsilon) = \sup\{|g(y)-g(z)| : |y-z| \leq \varepsilon\}$ be its
modulus of continuity, so $\omega(\varepsilon) \to 0$ as $\varepsilon \downarrow 0$.
For $0 < \varepsilon \leq 1$ and any $x$, using
$\int\rho_\varepsilon \, d\lambda = 1$ to write
$g(x) = \int \rho_\varepsilon(t) g(x)\, d\lambda(t)$ (legitimate because $g$ is
bounded and $\rho_\varepsilon$ is integrable),
\[
  \bigl|(\rho_\varepsilon * g)(x) - g(x)\bigr|
    = \Bigl|\int \rho_\varepsilon(t)\bigl(g(x-t)-g(x)\bigr) d\lambda(t)\Bigr|
    \;\leq\; \int \rho_\varepsilon(t)\,\bigl|g(x-t)-g(x)\bigr|\, d\lambda(t)
    \;\leq\; \omega(\varepsilon) ,
\]
the first inequality by Jensen's inequality (shared context) applied to the
probability measure $\mu_\varepsilon = \rho_\varepsilon\,d\lambda$ of Step 1 and
the convex function $s \mapsto |s|$, and the rest $\rho_\varepsilon(t) = 0$ unless $|t| \leq \varepsilon$, in which case
$|g(x-t)-g(x)| \leq \omega(\varepsilon)$, and again
$\int\rho_\varepsilon\,d\lambda = 1$.  Therefore
\begin{equation*}
  \|\rho_\varepsilon * g - g\|_{L^2(K,\lambda)}
    \;\leq\; \omega(\varepsilon)\, \lambda(K)^{1/2}
  \qquad (0 < \varepsilon \leq 1).
  \tag{2}
\end{equation*}

\emph{Step 3: approximation and conclusion.}
Let $\eta > 0$.  Put $f_1 = \mathbf{1}_{K_2} f$.  Then $f_1$ is measurable,
vanishes off the compact set $K_2$, and
$\int_{\mathbb{R}^n}|f_1|^2 d\lambda = \int_{K_2}|f|^2 d\lambda < \infty$ by
hypothesis, so $f_1 \in L^2(\mathbb{R}^n,\lambda)$.  Lebesgue measure is a regular
Borel measure, so by the density item of shared context there is
$g : \mathbb{R}^n \to \mathbb{R}$ continuous with compact support and
\begin{equation*}
  \|f_1 - g\|_{L^2(\mathbb{R}^n,\lambda)} \;\leq\; \eta/3 .
  \tag{3}
\end{equation*}

Next, for $0 < \varepsilon \leq 1$ and $x \in K$ we have
$(\rho_\varepsilon * f)(x) = (\rho_\varepsilon * f_1)(x)$: the integrand vanishes
unless $|t| \leq \varepsilon \leq 1$, and then $x - t \in K_1 \subseteq K_2$,
where $f_1 = f$.  Also $f = f_1$ on $K \subseteq K_2$.  Hence, pointwise on $K$,
\[
  \rho_\varepsilon * f - f
    = \rho_\varepsilon * (f_1 - g) \;+\; \bigl(\rho_\varepsilon * g - g\bigr)
      \;+\; \bigl(g - f_1\bigr) ,
\]
where we used linearity of the convolution in its second argument, legitimate
because $f_1 - g$ and $g$ are each in $L^2(K_1,\lambda)$ (the first by $(3)$, the
second because $g$ is bounded and $\lambda(K_1)<\infty$), so all three
convolutions are defined by Step 0.  Taking $L^2(K,\lambda)$ norms and using the
triangle inequality for the norm of $L^2(K,\lambda)$,
\[
  \|\rho_\varepsilon * f - f\|_{L^2(K,\lambda)}
    \;\leq\; \|\rho_\varepsilon * (f_1-g)\|_{L^2(K,\lambda)}
      + \|\rho_\varepsilon * g - g\|_{L^2(K,\lambda)}
      + \|g - f_1\|_{L^2(K,\lambda)} .
\]
By $(1)$ applied to $h = f_1 - g$, the first term is at most
$\|f_1-g\|_{L^2(K_1,\lambda)} \leq \|f_1-g\|_{L^2(\mathbb{R}^n,\lambda)} \leq \eta/3$
by $(3)$.  The third term is at most
$\|g-f_1\|_{L^2(\mathbb{R}^n,\lambda)} \leq \eta/3$, again by $(3)$.  By $(2)$ the
second term is at most $\omega(\varepsilon)\lambda(K)^{1/2}$, and since
$\omega(\varepsilon) \to 0$ there is $\varepsilon_0 \in (0,1]$ with
$\omega(\varepsilon)\lambda(K)^{1/2} \leq \eta/3$ for all
$0 < \varepsilon \leq \varepsilon_0$.  Altogether
\[
  \|\rho_\varepsilon * f - f\|_{L^2(K,\lambda)} \;\leq\; \eta
  \qquad\text{for all } 0 < \varepsilon \leq \varepsilon_0 .
\]
Since $\eta > 0$ was arbitrary, $\|\rho_\varepsilon * f - f\|_{L^2(K,\lambda)} \to 0$
as $\varepsilon \downarrow 0$ through positive values, which is the assertion.
\end{proof}
